% CV 1: Pour poste CDI en Intelligence Artificielle et Data Science
\documentclass[11pt,a4paper,sans]{moderncv}
\moderncvstyle{classic}
\moderncvcolor{blue}
\usepackage[utf8]{inputenc}
% \usepackage[scale=0.8, top=2cm, bottom=2cm, left=1.8cm, right=1.8cm]{geometry} % حاشیه‌های کاهش یافته
\usepackage[scale=0.8, top=1.2cm, bottom=1.2cm, left=1.2cm, right=1.2cm]{geometry}
% تنظیم عرض ستون Hints (عنوان‌ها) به صورت دستی
\setlength{\hintscolumnwidth}{3.8cm} % افزایش عرض به 4.5cm (می‌توانید این مقدار را تنظیم کنید)

\name{\cvfirstname}{{\cvlastname}}}
\setlength{\makecvtitlenamewidth}{10cm}
% \title{\bfseries Data Scientist \& Ingénieur IA}
\title{\bfseries Développeur Python}
% \address{\cvlocation}{}
% \address{\cvlocation}{}
\address{\cvlocation}{}
\phone[mobile]{\cvphone}
\email{\cvemail}
\social[linkedin]{\cvlinkedin}
\social[github]{\cvgithub}
\social[homepage]{\cvwebsite}
% 






% ── Personal data placeholders (overridden by ../shared/personal_data.tex) ────
\providecommand{\cvfirstname}{Firstname}
\providecommand{\cvlastname}{LASTNAME}
\providecommand{\cvemail}{user@example.com}
\providecommand{\cvphone}{+XX X XX XX XX XX}
\providecommand{\cvlocation}{City, Country}
\providecommand{\cvgithub}{github-username}
\providecommand{\cvlinkedin}{linkedin-username}
\providecommand{\cvwebsite}{website-url}
\input{../shared/personal_data}
\begin{document}
\makecvtitle

\vspace{-01.2cm}
\section{Profil}
\cvitem{}{
Développeur Python expérimenté, diplômé d'un Master en Informatique de l'Université de Montpellier.  
Expert en développement d'applications Python, automatisation de processus et analyse de données.  
Résultats concrets : optimisation de processus (+500\% vitesse, -99\% réduction de coûts, -97.7\% temps d'analyse). 
Disponible immédiatement.
}

\vspace{-0.3cm}
\section{Compétences}

% \cvitem{Data Visualization}{Matplotlib, Seaborn, Plotly, Tableau, Power BI}
% \cvitem{Automatisation}{Make (Integromat), Zapier, n8n, IFTTT \textit{(Low-Code/No-Code)}}

\cvitem{Data Science Python}{\textbf{Python}, Pandas, NumPy, Scikit-Learn, TensorFlow, LangChain, CrewAI}
\cvitem{Data \& Databases}{SQL, NoSQL, SQLAlchemy, MongoDB, Power BI, Matplotlib}
\cvitem{Outils \& DevOps}{Docker, Git, AWS, Flask, Streamlit, Agile SCRUM}

\vspace{-0.3cm}
\section{Expériences Professionnelles}

\cventry{Jan 2025 -- Juin 2025}{Développeur Python (Data \& IA)}{Stage -- toHero}{Montpellier}{}{
\begin{itemize}
\item \textbf{Automatisation IA:} Développement d'un système d'analyse automatisé avec validation par un expert métier : +500\% vitesse, -99\% coûts
\item \textbf{Automatisation Data:} Traitement automatisé de 18 000+ lignes Excel et 1 772 pages PDF via API OpenAI et Python : réduction de 97,7\% du temps d'analyse (88h → 2h)
\item \textbf{Tests multi-agents (CrewAI):} Génération de 476 scénarios de test fonctionnels pour une plateforme e-commerce, simulant des parcours utilisateurs variés
\end{itemize}
}

\cventry{2021 -- 2024}{Développeur Python \& Automatisation}{Compagnie Nationale du Pétrole Iranienne (NIOC)}{Iran}{}{
\begin{itemize}%
\item \textbf{Automatisation:} Développement de scripts Python pour automatiser les tâches réseau et analyser les données de performance, réduisant les interventions manuelles de 20\%.
\item \textbf{Gestion d'infrastructure:} Surveillance et sécurisation des équipements Cisco pour un réseau de +1500 utilisateurs.
\item \textbf{Compétences :} Python avancé, visualisation de données (rapports interactifs avec Power BI), résolution de problèmes complexes, pensée algorithmique.
\end{itemize}
}

% \cventry{2017 -- 2019}{Ingénieur Réseaux (VOIP)}{NIOC}{Iran}{}{
% \begin{itemize}%
%     \item \textbf{Analyse de données de performance :} Collaboration avec une équipe de 10 techniciens pour la configuration, maintenance et analyse des données de performance des systèmes VOIP pour +1500 utilisateurs.
% \item \textbf{Support technique:} Gestion des serveurs et routeurs, développant une expertise en infrastructure IT.
% \end{itemize}
% }

% \cventry{2010 -- 2013}{Inspecteur QA Projets Télécommunications}{NIOC}{Iran}{}{
% {\footnotesize
% \begin{itemize}%
% \item \textbf{Gestion de projet:} Supervision de l'installation d'un réseau TETRA sur 78 sites.
% \item \textbf{Contrôle qualité:} Vérification de la conformité des équipements.
% \end{itemize}
% }
% }

\vspace{-0.3cm}
\section{Projets}
\cventry{2025 -- 2026}{Développement et Déploiement d'une Application d'Analyse de Données}{Projet d'Université (DU Python)}{}{}{
    \begin{itemize}
        \item \textbf{Développement Python:} Mise en œuvre de scripts Python pour le nettoyage, l'intégration et l'analyse de données d'assurance complexes (\textbf{922 000+ enregistrements historiques}).        
        \item \textbf{Modélisation \& Algorithmes:} Application d'analyses économétriques et d'algorithmes de Machine Learning pour l'identification de relations et d'anomalies (approche algorithmique).
        \item \textbf{Mise en Production (Streamlit):} Conception et déploiement d'une application web interactive de DataViz via Streamlit, assurant une interface utilisateur fonctionnelle.
    \end{itemize}
\scriptsize{\textit{Technologies: Python, programmation orientée objet (POO), Machine Learning, Streamlit, Économétrie, Pandas}}
}
% \cventry{2023}{Comparaison de Modèles de Deep Learning pour la Prédiction Boursière}{Projet Freelance}{}{}{
% \begin{itemize}
% \item Analyse comparative de 4 architectures de deep learning (LSTM, BiLSTM, CNN-BiLSTM, R-CNN-BiLSTM) sur un portefeuille de 30 actions.
% \item Prétraitement des données financières et feature engineering en Python.
% \item Identification du modèle optimal (CNN-BiLSTM) pour maximiser la précision des prévisions.
% \end{itemize}
% \scriptsize{\textit{Technologies: Python, Keras, TensorFlow, Pandas, NumPy}}
% }

\cventry{2023}{Prédiction Boursière par Deep Learning}{Projet Freelance}{}{}{
    \begin{itemize}
        \item Comparaison de 3 modèles (LSTM, BiLSTM, CNN-BiLSTM) - optimisation précision
    \end{itemize}
\scriptsize{\textit{Technologies: Python, Keras, TensorFlow, Pandas, NumPy}}
}


% \cventry{2022}{Analyse de Données Financières et Web Scraping}{Projet Personnel}{}{}{
% \begin{itemize}
% \item Collecte automatisée et nettoyage de données financières via des techniques de web scraping (Beautiful Soup, Scrapy).
% \item Exploration et analyse de données pour identifier des tendances de marché.
% \item Entraînement de modèles de machine learning (régression, classification) pour des analyses prédictives.
% \end{itemize}
% \scriptsize{\textit{Technologies: Python, Beautiful Soup, Pandas, Scikit-Learn}}
% }


\cventry{2022}{Web Scraping \& Analyse Financière}{Projet Personnel}{}{}{
    \begin{itemize}
        \item Collecte et analyse de données financières via web scraping + modèles ML pour prédictions
    \end{itemize}
\scriptsize{\textit{Technologies: Python, Beautiful Soup, Pandas, Scikit-Learn}}
}


% \cventry{2020}{Application Web de Gestion de Données Culturelles}{Projet Univ. Montpellier}{}{}{
% \begin{itemize}
% \item Conception et développement full-stack d'une application de gestion et visualisation de données.
% \item Mise en place d'une API RESTful avec Spring Boot et base de données MySQL.
% \end{itemize}
% \scriptsize{\textit{Technologies: Java, Spring Boot, MySQL, JavaScript}}
% }

% \cventry{2020}{E-Commerce SPA: Panier \& Authentification (Angular + Node.js)}{Projet Univ. Montpellier}{}{}{
% \begin{itemize}
% \item Développement d’une application monopage (SPA) de panier e-commerce avec authentification JWT, gestion des commandes utilisateur et recherche produit.
% \item Implémentation backend Node.js : tri alphabétique serveur, upload photo (Multer), hashing mot de passe (bcryptjs), guards Angular, MVC, gestion d’erreurs.
% \item Frontend Angular 10 : interface responsive, barre de recherche, ajout/suppression panier, validation formulaires, affichage détails produit/commande.
% \item Base MongoDB conteneurisée avec Docker.
% \end{itemize}
% \scriptsize{\textit{Technologies: Node.js, Express, Angular 10, MongoDB, JWT, Docker, Git, REST API, MVC}}
% }

% \cventry{2019}{Enrichissement d'une Base de Connaissances pour Fact-Checking}{Projet Univ. Montpellier}{}{}{
% \begin{itemize}
% \item Development d'un pipeline de collecte et de traitement de données pour +10k articles.
% \item Nettoyage et structuration des données pour alimenter un système de recherche d'information.
% \end{itemize}
% \scriptsize{\textit{Technologies: Python, Beautiful Soup, Pandas}}
% }

% \cventry{2019}{Plateforme d'Exploration de Données Financières en Temps Réel}{Projet Univ. Montpellier}{}{}{
% \begin{itemize}
% \item Création d'une application web interactive pour l'extraction et la visualisation de données financières.
% \item Mise en œuvre d'une recherche instantanée et de graphiques dynamiques.
% \end{itemize}
% \scriptsize{\textit{Technologies: Python, Flask, MongoDB, Plotly}}
% }

\vspace{-0.3cm}
\section{Formations}
\cventry{2025 -- 2026}{Diplôme d'Université Big Data, Data Science et Analyse des Risques sous Python (DU, obtenu)}{\textbf{Faculté d'Économie}, Université de Montpellier, France}{}{}{}
\cventry{2024 -- 2025}{Master en Informatique (Bac+5)}{Université de Montpellier, France}{}{}{}
% \cventry{2015}{Licence en Génie Électronique}{Université Jahad-e-Daneshgahi Khouzestan, Iran}{}{GPA: 17.94/20}{}

% \section{Certifications}
% \cvitem{Certifications}{Multi-Agent Systems (CrewAI), Prompt Engineering, Machine Learning Bootcamp}

% \cvitem{2025}{Multi-Agent Systems (CrewAI)  }
% \cvitem{2025}{ChatGPT Prompt Engineering (DeepLearning.AI)  }
% \cvitem{2023}{Bootcamp Apprentissage Automatique (OnAcademy)  }
% \cvitem{2022}{Bootcamp Analyse Financière (RoseAcademy)  }

% \section{Langues}
% \cvitem{Français}{Avancé}
% \cvitem{Anglais}{Niveau Opérationnel Professionnel}

% \section{Références}
% \cvitem{}{Références disponibles sur demande}

\vspace{-0.3cm}
\section{Langues \& Certifications}
\cvitem{Langues}{Français (Avancé), Anglais (Opérationnel Professionnel), Persan (Natif)}
\cvitem{Certifications}{Multi-Agent Systems (CrewAI), Prompt Engineering, Machine Learning Bootcamp}

\end{document}






